\documentclass[11pt,a4paper]{article}
% =====================================================================
%  Gemeinsame Vorlage der Arbeitsblätter "Rekursive Datenstrukturen"
%  – Informatik 12 (gleiches Layout wie informatik/13/formale-sprachen/
%  arbeitsblaetter/ab-vorlage.tex, dazu Listings und Objektdiagramme)
%  Einbinden in jedem Blatt direkt nach \documentclass: % =====================================================================
%  Gemeinsame Vorlage der Arbeitsblätter "Rekursive Datenstrukturen"
%  – Informatik 12 (gleiches Layout wie informatik/13/formale-sprachen/
%  arbeitsblaetter/ab-vorlage.tex, dazu Listings und Objektdiagramme)
%  Einbinden in jedem Blatt direkt nach \documentclass: % =====================================================================
%  Gemeinsame Vorlage der Arbeitsblätter "Rekursive Datenstrukturen"
%  – Informatik 12 (gleiches Layout wie informatik/13/formale-sprachen/
%  arbeitsblaetter/ab-vorlage.tex, dazu Listings und Objektdiagramme)
%  Einbinden in jedem Blatt direkt nach \documentclass: \input{ab-vorlage}
%  Übersetzen mit XeLaTeX, LuaLaTeX, pdfLaTeX oder Tectonic.
%
%  Blatt:   xelatex ab-1-2-interfaces.tex
%  Lösung:  xelatex ab-1-2-interfaces-lsg.tex
%           (Datei enthält nur: \def\mitloesung{}\input{ab-1-2-interfaces};
%            füllt alle Lücken rot aus; siehe \lsg, \lsglinien, \lsgoder,
%            \lsgtext, lsgbild, \kreuz)
% =====================================================================
\usepackage[a4paper,left=16mm,right=16mm,top=23mm,bottom=13mm,headheight=12mm,headsep=3mm,footskip=6mm]{geometry}
\usepackage{iftex}
\ifPDFTeX
  \usepackage[T1]{fontenc}
  \usepackage[utf8]{inputenc}
  \usepackage{helvet}
  \usepackage[scaled=0.86]{beramono}
  \renewcommand{\familydefault}{\sfdefault}
\else
  \usepackage{fontspec}
  \setmainfont{texgyreheros}[Extension=.otf,UprightFont=*-regular,BoldFont=*-bold,ItalicFont=*-italic,BoldItalicFont=*-bolditalic]
  \setmonofont{DejaVuSansMono}[Extension=.ttf,UprightFont=*,BoldFont=*-Bold,ItalicFont=*-Oblique,BoldItalicFont=*-BoldOblique,Scale=0.86]
\fi
\usepackage[ngerman,shorthands=off]{babel}
\usepackage{amssymb}
\usepackage{xcolor,graphicx,tabularx,array,enumitem,fancyhdr,tikz,listings,multicol,calc,etoolbox,pifont}
\usepackage[most]{tcolorbox}
\usetikzlibrary{arrows.meta,positioning,calc,decorations.pathreplacing,shapes.multipart}
\usepackage[hidelinks]{hyperref}
\setlength{\parindent}{0pt}
\setlength{\parskip}{3pt}
\setlist{nosep,leftmargin=*}

% ---------- Lösungsschalter ----------
\newif\ifloesung
\ifdefined\mitloesung\loesungtrue\fi
\newcommand{\lsgfarbe}{\color{red!75!black}}

% ---------- Kopf- und Fußzeile (einheitliche Form) ----------
\newcommand{\lehrkraft}{Hau}
\newcommand{\themenbereich}{1. Rekursive Datenstrukturen}
\newcommand{\abnummer}{}
\pagestyle{fancy}\fancyhf{}
\renewcommand{\headrulewidth}{0.4pt}
\fancyhead[L]{\small Klasse 12\\Datum:}
\fancyhead[C]{\small Themenbereich:\\\themenbereich}
\fancyhead[R]{\small \lehrkraft\ifloesung\\{\lsgfarbe\bfseries Lösung}\fi}
\fancyfoot[C]{\scriptsize Informatik 12\,·\,Blatt \abnummer\ifloesung\ (Lösung)\fi\,·\,Seite \thepage}
\newcommand{\abtitel}[2]{\renewcommand{\abnummer}{#1}\begin{center}\Large\bfseries #1\quad\underline{#2}\end{center}\vspace{-2mm}}

% ---------- Boxen ----------
\newenvironment{aufgabe}[2][]{\begin{tcolorbox}[enhanced,breakable,colback=white,colframe=black,boxrule=0.7pt,arc=3mm,left=2mm,right=2mm,top=1mm,bottom=1.2mm,before skip=2mm,after skip=2mm]\textbf{Aufgabe #2}\ifstrempty{#1}{}{ \textit{#1}}\par\vspace{0.3mm}}{\end{tcolorbox}}
\newtcolorbox{merke}[1][Merke]{enhanced,breakable,colback=green!5,colframe=green!40!black,boxrule=0.8pt,arc=2mm,left=2mm,right=2mm,top=1mm,bottom=1.2mm,before skip=2mm,after skip=2mm,title={#1},fonttitle=\bfseries,coltitle=white,toptitle=0.3mm,bottomtitle=0.3mm}
\newtcolorbox{hinweis}{enhanced,breakable,colback=yellow!12,colframe=orange!70!black,boxrule=0.6pt,arc=2mm,left=2mm,right=2mm,top=1mm,bottom=1mm,before skip=2mm,after skip=2mm}
\newcommand{\web}[1]{{\small\color{gray}\textit{Website: #1}}}

% ---------- Lücken, Linien, Karos ----------
\newcommand{\luecke}[1][3.5cm]{\rule[-0.5ex]{#1}{0.4pt}}
\newcommand{\linien}[1]{\par\vspace{1.5mm}\foreach \i in {1,...,#1}{\noindent\rule[0pt]{\linewidth}{0.3pt}\par\vspace{5.5mm}}\vspace{-3mm}}
\newcommand{\kariert}[1]{\par\vspace{1.5mm}\noindent\begin{tikzpicture}\draw[step=5mm,gray!55,thin] (0,0) grid (\linewidth-1pt,#1*5mm);\end{tikzpicture}\par}
\newcommand{\karobox}[2]{\begin{tikzpicture}\draw[step=5mm,gray!55,thin] (0,0) grid (#1,#2);\end{tikzpicture}}
\newcommand{\code}[1]{\texttt{#1}}
% Schritt-Nummern wie im Code und auf der Website: \nr{1} = ①
\newcommand{\nr}[1]{\ding{\numexpr171+#1\relax}}

% Lücke mit Lösung: \lsg[Mindestbreite]{Lösungstext}
%  Die Lösung steht (ohne eigene Höhe/Tiefe) auf derselben Linie wie im Blatt: Das Layout verschiebt sich nicht,
%  solange der Text nicht breiter als die Mindestbreite ist.
\newlength{\lsgbreite}
\newcommand{\lsg}[2][3.5cm]{\leavevmode\ifloesung\settowidth{\lsgbreite}{\,#2\,}\ifdim\lsgbreite<#1\setlength{\lsgbreite}{#1}\fi\rlap{\smash{\raisebox{0.8pt}{\makebox[\lsgbreite][l]{\lsgfarbe\,#2}}}}\luecke[\lsgbreite]\else\luecke[#1]\fi}
% Schreiblinien mit Lösung: \lsglinien{Anzahl}{Lösungstext} (gleiche Höhe wie \linien:
%  je Linie \baselineskip + 5,5 mm + \parskip, ausgemessen)
\newlength{\lsglinienhoehe}
\newcommand{\lsglinien}[2]{\ifloesung\par\vspace{1.5mm}\setlength{\lsglinienhoehe}{\dimexpr #1\dimexpr\baselineskip+5.5mm+\parskip\relax-16.25pt\relax}\noindent\begin{minipage}[t][\lsglinienhoehe][t]{\linewidth}\lsgfarbe #2\end{minipage}\par\vspace{-1.5mm}\else\linien{#1}\fi}
% Beliebiger Inhalt: \lsgoder{nur in der Lösung}{nur im Blatt}
\newcommand{\lsgoder}[2]{\ifloesung\leavevmode#1\else#2\fi}
% Text nur in der Lösung sichtbar, belegt im Blatt aber denselben Platz (ohne Text im PDF)
\newcommand{\lsgtext}[1]{\leavevmode\ifloesung{\lsgfarbe#1}\else\phantom{#1}\fi}
% Ankreuzfeld: \kreuz{1} = angekreuzt (nur in der Lösung), \kreuz{0} = leer
\newcommand{\kreuz}[1]{\leavevmode\ifloesung\ifnum#1=1 {\lsgfarbe$\boxtimes$}\else$\square$\fi\else$\square$\fi}
% Tabellenzeile mit fester Höhe für Lösungseinträge (2 Zeilen \small)
\newcommand{\zeile}{\rule[-6mm]{0pt}{9.5mm}}

% ---------- Verkettete Strukturen zeichnen ----------
% Objekt mit Kopf und Attributen:  \node[obj] (k1) {k1: Knoten \nodepart{second} daten = Hans\\ nachfolger};
\tikzset{
  obj/.style={draw,rectangle split,rectangle split parts=2,rectangle split part fill={orange!25,orange!8},
              rounded corners=2pt,font=\scriptsize,inner sep=2pt,align=left},
  qobj/.style={obj,rectangle split part fill={teal!25,teal!8}},
  aobj/.style={obj,rectangle split part fill={gray!30,gray!8}},
  pobj/.style={draw,rounded corners=2pt,fill=gray!8,font=\scriptsize,inner sep=2pt,align=left},
  klasse/.style={draw,rectangle split,rectangle split parts=3,font=\scriptsize,align=left,fill=white},
  zeiger/.style={thick,-{Stealth[length=2.2mm]}},
  erbt/.style={thick,-{Triangle[open,length=2.6mm,width=2.4mm]}},
  implementiert/.style={thick,densely dashed,-{Triangle[open,length=2.6mm,width=2.4mm]}},
  lsgzeiger/.style={zeiger,red!75!black},
  nullref/.style={font=\scriptsize\ttfamily,inner sep=1pt},
  schritt/.style={circle,draw,fill=white,inner sep=0.6pt,font=\scriptsize\bfseries},
  lsgschritt/.style={schritt,draw=red!75!black,text=red!75!black},
  % Sequenzdiagramm (UML 2): Kopf eckig, synchron = gefüllte Spitze, asynchron = offene Spitze,
  %  Rückgabe und Erzeugung gestrichelt mit offener Spitze (kein hohles Dreieck: das wäre Vererbung)
  lebenslinie/.style={draw,sharp corners,fill=white,font=\scriptsize,inner sep=2pt},
  leben/.style={densely dashed,gray},
  aufruf/.style={thick,-{Stealth[length=2mm]}},
  asynchron/.style={thick,-{Straight Barb[length=2mm,width=2.4mm]}},
  rueckgabe/.style={thick,densely dashed,-{Straight Barb[length=2mm,width=2.4mm]}},
  erzeugung/.style={rueckgabe},
}
% Akteur (Strichmännchen) mit Namen darunter: \akteur{x}{y oben}{Name}; ca. 0,9 cm hoch
\newcommand{\akteur}[3]{\begin{scope}[shift={(#1,#2)},line width=0.6pt]
  \draw (0,-0.09cm) circle (0.09cm);
  \draw (0,-0.18cm) -- (0,-0.42cm) (-0.16cm,-0.26cm) -- (0.16cm,-0.26cm)
        (-0.13cm,-0.62cm) -- (0,-0.42cm) -- (0.13cm,-0.62cm);
  \node[font=\scriptsize,inner sep=1pt,anchor=north] at (0,-0.63cm) {#3};
\end{scope}}
% Aktivierungsbalken auf der Lebenslinie bei x von y1 bis y2: \aktiv{x}{y1}{y2}
\newcommand{\aktiv}[3]{\draw[fill=gray!15] ($(#1,#2)+(-0.09cm,0)$) rectangle ($(#1,#3)+(0.09cm,0)$);}
% Lösungszeichnung: im Blatt unsichtbar, belegt aber denselben Platz
\ifloesung
  \tikzset{lsgbild/.style={red!75!black}}
\else
  \tikzset{lsgbild/.style={draw opacity=0,fill opacity=0,text opacity=0}}
\fi

% ---------- Java-Listings ----------
\lstset{language=Java,basicstyle=\ttfamily\small,keywordstyle=\bfseries,commentstyle=\color{gray!80!black},stringstyle=\color{black},showstringspaces=false,columns=flexible,keepspaces=true,tabsize=3,frame=single,framerule=0.4pt,rulecolor=\color{black},xleftmargin=2mm,framexleftmargin=1mm,aboveskip=2mm,belowskip=2mm,breaklines=true,escapechar=§,
 morekeywords={interface,implements,abstract,extends},
 literate={ä}{{\"a}}1 {ö}{{\"o}}1 {ü}{{\"u}}1 {Ä}{{\"A}}1 {Ö}{{\"O}}1 {Ü}{{\"U}}1 {ß}{{\ss}}1 {–}{{--}}1 {„}{{\glqq}}1 {“}{{\grqq}}1 {→}{{$\to$}}1}
\lstnewenvironment{java}[1][]{\lstset{#1}}{}

%  Übersetzen mit XeLaTeX, LuaLaTeX, pdfLaTeX oder Tectonic.
%
%  Blatt:   xelatex ab-1-2-interfaces.tex
%  Lösung:  xelatex ab-1-2-interfaces-lsg.tex
%           (Datei enthält nur: \def\mitloesung{}\documentclass[11pt,a4paper]{article}
\input{ab-vorlage}
\begin{document}
\abtitel{1.2}{Interfaces: ein Vertrag für alle Daten}

Am Bahnhof warten Taxis in einer Schlange: Wer zuerst ankommt, fährt zuerst. Die Stadt möchte dafür die \code{Warteschlange} aus 1.1 benutzen. Doch die kann bisher nur Patienten verwalten. Eine Kopie für Taxis (\code{TaxiKnoten}, \code{TaxiWarteschlange}) würde funktionieren, aber jede Verbesserung müsste man dann doppelt einbauen. \web{Station 1.2, Projekt Taxistand}

\begin{aufgabe}[Was muss die Warteschlange über ihre Daten wissen?]{1}
Untersuche im Projekt die Methoden der Warteschlange. Kreuze an, ob die Methode die Daten nur weiterreicht oder eine Methode der Daten aufruft. Notiere diese gegebenenfalls.\\[1mm]
\small\renewcommand{\arraystretch}{1.2}
\begin{tabularx}{\linewidth}{|l|c|c|X|}\hline
\textbf{Methode} & \textbf{reicht nur weiter} & \textbf{ruft eine Methode der Daten auf} & \textbf{welche?}\\\hline
\code{add(Patient p)} & \kreuz{1} & \kreuz{0} & \\\hline
\code{poll()}, \code{peek()} & \kreuz{1} & \kreuz{0} & \\\hline
\code{ausgeben()} & \kreuz{0} & \kreuz{1} & \lsgtext{\code{getName()}: der Name zum Ausgeben}\\\hline
\end{tabularx}\\[1.5mm]
\normalsize
\textbf{Folgerung:} Von ihren Daten muss die Warteschlange nur verlangen, dass sie einen \lsg[4.6cm]{Namen (Schlüsselwert) nennen} können. Genau das legen wir in einem Vertrag fest: im \textbf{Interface} \code{DatenElement}.
\end{aufgabe}

\begin{minipage}[t]{0.47\linewidth}
\vspace{0pt}
\begin{java}[basicstyle=\ttfamily\footnotesize]
public §\lsg[1.9cm]{interface}§ DatenElement {
   String gibSchluesselwert();
}

public class Taxi §\lsg[3.2cm]{implements DatenElement}§ {
   private int nummer;
   private String fahrer;
   ...
   public String gibSchluesselwert() {
      return §\lsg[2.6cm]{"Taxi " + nummer}§;
   }
}
\end{java}
\end{minipage}\hfill
\begin{minipage}[t]{0.505\linewidth}
\vspace{0pt}
\begin{merke}[Interface]
\small
\begin{itemize}
\item Ein Interface ist ein \lsg[1.7cm]{Vertrag}. Es legt nur \lsg[2.9cm]{Methodenköpfe} fest, keine Rümpfe.
\item Eine Klasse, die es implementiert (\code{implements}), muss \lsg[1.3cm]{alle} seine Methoden bereitstellen. Das prüft der \lsg[2cm]{Compiler}.
\item Ein Interface ist auch ein \lsg[1.9cm]{Datentyp}: Eine Variable vom Typ \code{DatenElement} kann Objekte jeder Klasse aufnehmen, die es implementiert. Objekte vom Interface selbst gibt es \lsg[1.2cm]{nicht}.
\end{itemize}
\end{merke}
\end{minipage}

\begin{aufgabe}[Klassendiagramm]{2}
\textbf{a)} Zeichne alle Beziehungen ein: Assoziationen mit Namen und Multiplizitäten, dazu \code{implements} als gestrichelten Pfeil mit leerem Dreieck zum Interface.\\
\textbf{b)} Umrande die Klassen der \textbf{Struktur} blau und die Klassen des \textbf{Inhalts} grün.\\[1mm]
\begin{tikzpicture}
\node[klasse,fill=teal!6,anchor=north west] (W) at (0,0)
  {\textbf{Warteschlange}\nodepart{second}- anfang: Knoten\\- ende: Knoten
   \nodepart{third}+ add(DatenElement d): boolean\\+ poll(): DatenElement\\+ ausgeben()};
\node[klasse,fill=orange!8,anchor=north west] (K) at (6.0,0)
  {\textbf{Knoten}\nodepart{second}- daten: DatenElement\\- nachfolger: Knoten
   \nodepart{third}+ getDaten(): DatenElement\\+ getNachfolger(): Knoten};
\node[klasse,fill=green!6,anchor=north,rectangle split parts=2] (D) at (13.1,0)
  {\textit{\guillemotleft interface\guillemotright}\\\textbf{DatenElement}\nodepart{second}+ gibSchluesselwert(): String};
\node[klasse,fill=gray!8,anchor=north] (P) at (11.2,-3.0)
  {\textbf{Patient}\nodepart{second}- name: String
   \nodepart{third}+ gibSchluesselwert(): String};
\node[klasse,fill=gray!8,anchor=north] (T) at (15.15,-3.0)
  {\textbf{Taxi}\nodepart{second}- nummer: int\\- fahrer: String
   \nodepart{third}+ gibSchluesselwert(): String};
\begin{scope}[font=\scriptsize]
  \draw[zeiger,lsgbild] ([yshift=-4mm]W.north east) -- ([yshift=-4mm]W.north east -| K.west)
     node[pos=0.3,above] {\lsgtext{anfang}} node[pos=0.86,below] {\lsgtext{0..1}};
  \draw[zeiger,lsgbild] ([yshift=-13mm]W.north east) -- ([yshift=-13mm]W.north east -| K.west)
     node[pos=0.3,below] {\lsgtext{ende}} node[pos=0.86,above] {\lsgtext{0..1}};
  \draw[zeiger,lsgbild] ([yshift=-5mm]K.north east) -- ([yshift=-5mm]K.north east -| D.west)
     node[pos=0.6,below] {\lsgtext{daten}} node[pos=0.88,above] {\lsgtext{1}};
  \draw[zeiger,lsgbild] ([xshift=-5mm]K.north east) -- ++(0,4mm) -| ([xshift=5mm]K.north west);
  \node[left] at ($(K.north west)+(4.5mm,3.2mm)$) {\lsgtext{nachfolger\quad 0..1}};
  \draw[implementiert,lsgbild] ([xshift=10mm]P.north) -- ([xshift=10mm]P.north |- D.south);
  \draw[implementiert,lsgbild] ([xshift=-9.5mm]T.north) -- ([xshift=-9.5mm]T.north |- D.south);
  \draw[lsgbild,blue!70!black,thick,rounded corners] ($(W.north west)+(-1.5mm,7.5mm)$) rectangle ($(K.east |- W.south)+(1.5mm,-1.5mm)$);
  % Inhalt: oben nur um das Interface, unten um Patient und Taxi (unter dem blauen Rahmen)
  \coordinate (gl) at ($(K.east)+(2.75mm,0)$);
  \coordinate (gt) at ($(D.north)+(0,2mm)$);
  \coordinate (gr) at ($(T.east)+(1.5mm,0)$);
  \coordinate (gb) at ($(T.south)+(0,-1.5mm)$);
  \coordinate (gp) at ($(P.west)+(-1.5mm,0)$);
  \coordinate (gm) at ($(W.south)+(0,-3mm)$);
  \draw[lsgbild,green!45!black,thick,rounded corners] (gl |- gt) -- (gr |- gt) -- (gr |- gb) -- (gp |- gb) -- (gp |- gm) -- (gl |- gm) -- cycle;
\end{scope}
\end{tikzpicture}\\[1mm]
\textbf{c)} Durch die Trennung von Struktur und Inhalt erledigt jede Klasse genau eine Aufgabe (hohe \lsg[2.1cm]{Kohäsion}). Struktur und Inhalt hängen nur über das Interface zusammen (lose \lsg[2.1cm]{Kopplung}). So kann die Warteschlange ohne jede Änderung \lsg[4.5cm]{jedes neue DatenElement} aufnehmen.
\end{aufgabe}

\newpage
\begin{aufgabe}[Projekt Taxistand, TODO 3a bis 3d]{3}
Baue das Projekt so um, dass Warteschlange und Knoten mit jedem \code{DatenElement} arbeiten. Ergänze die Tabelle und übertrage die Änderungen ins Projekt (Reihenfolge wie die TODOs).\\[1mm]
\small\renewcommand{\arraystretch}{1.25}
\begin{tabularx}{\linewidth}{|c|l|X|}\hline
 & \textbf{bisher} & \textbf{neu}\\\hline
3a & \code{DatenElement}: (leer) & \lsgtext{\code{String gibSchluesselwert();}}\\\hline
3b & \code{public class Patient \{} & \lsgtext{\code{public class Patient implements DatenElement \{}}\\
   & & \lsgtext{\code{+ public String gibSchluesselwert() \{ return name; \}}}\\\hline
3c & \code{private Patient daten;} & \lsgtext{\code{private DatenElement daten;}} \ (ebenso Konstruktor, \code{getDaten})\\\hline
3d & \code{public boolean add(Patient p)} & \lsgtext{\code{public boolean add(DatenElement daten)}}\\
   & \code{public Patient poll()} & \lsgtext{\code{public DatenElement poll()}} \ (ebenso \code{peek})\\
   & \code{aktuell.getDaten().getName()} & \lsgtext{\code{aktuell.getDaten().gibSchluesselwert()}}\\\hline
\end{tabularx}\\[1.5mm]
\normalsize
\textbf{b)} Nach dem Umbau von \code{Knoten} meldet der Compiler einen Fehler in \code{ausgeben()}. Erkläre!\\
\lsg[\linewidth]{Ein DatenElement hat nur die Methode gibSchluesselwert(), kein getName().}\\[1mm]
\lsg[\linewidth]{Die Warteschlange darf nur benutzen, was im Vertrag steht.}
\end{aufgabe}

\begin{aufgabe}[Projekt Taxistand, TODO 4]{4}
\textbf{a)} Ergänze den Code von \code{Taxi} oben auf Seite 1, übertrage ihn ins Projekt und kommentiere im Hauptprogramm die Taxi-Zeilen ein.\\
\textbf{b)} Lösche testweise die Methode \code{gibSchluesselwert()} aus \code{Taxi}. Notiere die Meldung des Compilers und erkläre sie mit dem Wort \textit{Vertrag}.\\
\lsg[\linewidth]{„Die Klasse Taxi muss noch folgende Methoden des Interfaces DatenElement implementieren: …“}\\[1mm]
\lsg[\linewidth]{Taxi verspricht mit implements, den Vertrag zu erfüllen, hält ihn aber nicht ein.}
\end{aufgabe}

\begin{aufgabe}[Datentypen]{5}
\code{taxistand} ist eine \code{Warteschlange}. Kreuze an, welche Zeilen der Compiler akzeptiert, und begründe kurz.\\[1mm]
\small\renewcommand{\arraystretch}{1.25}
\begin{tabularx}{\linewidth}{|l|c|c|X|}\hline
\textbf{Zeile} & \textbf{ok} & \textbf{Fehler} & \textbf{Begründung}\\\hline
\code{DatenElement d = new Taxi(5, "Lea");} & \kreuz{1} & \kreuz{0} & \lsgtext{Taxi implementiert DatenElement.}\\\hline
\code{DatenElement d = new DatenElement();} & \kreuz{0} & \kreuz{1} & \lsgtext{Von einem Interface gibt es keine Objekte.}\\\hline
\code{taxistand.add(new Patient("Ali", "Husten"));} & \kreuz{1} & \kreuz{0} & \lsgtext{Auch Patient ist ein DatenElement.}\\\hline
\code{Taxi t = taxistand.poll();} & \kreuz{0} & \kreuz{1} & \lsgtext{poll() liefert nur ein DatenElement.}\\\hline
\code{println(taxistand.poll().gibSchluesselwert());} & \kreuz{1} & \kreuz{0} & \lsgtext{Diese Methode steht im Vertrag.}\\\hline
\code{println(taxistand.poll().getFahrer());} & \kreuz{0} & \kreuz{1} & \lsgtext{getFahrer() steht nicht im Vertrag.}\\\hline
\end{tabularx}
\normalsize
\end{aufgabe}

\begin{hinweis}\small
\textbf{Es geht auch anders -- besprecht mündlich, warum man es so machen kann:}
\begin{itemize}
\item Jede Klasse ist ein \code{Object}. Der Knoten könnte also einfach \code{Object daten} speichern. Was ginge dann in \code{ausgeben()} nicht mehr?
\item In einer älteren Fassung stand im Vertrag nur \code{String toString();}. Warum erzwingt dieser Vertrag nichts? (Tipp: Was erbt jede Klasse von \code{Object}?)
\item Statt des Interfaces gäbe es eine Oberklasse \code{Daten}, von der Patient und Taxi erben. Was spricht dagegen, wenn \code{Patient} schon von \code{Person} erbt?
\item Du kennst schon ein Interface: \code{Wartezimmer} aus 1.1. Was legte es fest -- und wozu?
\end{itemize}
\end{hinweis}

\end{document}
;
%            füllt alle Lücken rot aus; siehe \lsg, \lsglinien, \lsgoder,
%            \lsgtext, lsgbild, \kreuz)
% =====================================================================
\usepackage[a4paper,left=16mm,right=16mm,top=23mm,bottom=13mm,headheight=12mm,headsep=3mm,footskip=6mm]{geometry}
\usepackage{iftex}
\ifPDFTeX
  \usepackage[T1]{fontenc}
  \usepackage[utf8]{inputenc}
  \usepackage{helvet}
  \usepackage[scaled=0.86]{beramono}
  \renewcommand{\familydefault}{\sfdefault}
\else
  \usepackage{fontspec}
  \setmainfont{texgyreheros}[Extension=.otf,UprightFont=*-regular,BoldFont=*-bold,ItalicFont=*-italic,BoldItalicFont=*-bolditalic]
  \setmonofont{DejaVuSansMono}[Extension=.ttf,UprightFont=*,BoldFont=*-Bold,ItalicFont=*-Oblique,BoldItalicFont=*-BoldOblique,Scale=0.86]
\fi
\usepackage[ngerman,shorthands=off]{babel}
\usepackage{amssymb}
\usepackage{xcolor,graphicx,tabularx,array,enumitem,fancyhdr,tikz,listings,multicol,calc,etoolbox,pifont}
\usepackage[most]{tcolorbox}
\usetikzlibrary{arrows.meta,positioning,calc,decorations.pathreplacing,shapes.multipart}
\usepackage[hidelinks]{hyperref}
\setlength{\parindent}{0pt}
\setlength{\parskip}{3pt}
\setlist{nosep,leftmargin=*}

% ---------- Lösungsschalter ----------
\newif\ifloesung
\ifdefined\mitloesung\loesungtrue\fi
\newcommand{\lsgfarbe}{\color{red!75!black}}

% ---------- Kopf- und Fußzeile (einheitliche Form) ----------
\newcommand{\lehrkraft}{Hau}
\newcommand{\themenbereich}{1. Rekursive Datenstrukturen}
\newcommand{\abnummer}{}
\pagestyle{fancy}\fancyhf{}
\renewcommand{\headrulewidth}{0.4pt}
\fancyhead[L]{\small Klasse 12\\Datum:}
\fancyhead[C]{\small Themenbereich:\\\themenbereich}
\fancyhead[R]{\small \lehrkraft\ifloesung\\{\lsgfarbe\bfseries Lösung}\fi}
\fancyfoot[C]{\scriptsize Informatik 12\,·\,Blatt \abnummer\ifloesung\ (Lösung)\fi\,·\,Seite \thepage}
\newcommand{\abtitel}[2]{\renewcommand{\abnummer}{#1}\begin{center}\Large\bfseries #1\quad\underline{#2}\end{center}\vspace{-2mm}}

% ---------- Boxen ----------
\newenvironment{aufgabe}[2][]{\begin{tcolorbox}[enhanced,breakable,colback=white,colframe=black,boxrule=0.7pt,arc=3mm,left=2mm,right=2mm,top=1mm,bottom=1.2mm,before skip=2mm,after skip=2mm]\textbf{Aufgabe #2}\ifstrempty{#1}{}{ \textit{#1}}\par\vspace{0.3mm}}{\end{tcolorbox}}
\newtcolorbox{merke}[1][Merke]{enhanced,breakable,colback=green!5,colframe=green!40!black,boxrule=0.8pt,arc=2mm,left=2mm,right=2mm,top=1mm,bottom=1.2mm,before skip=2mm,after skip=2mm,title={#1},fonttitle=\bfseries,coltitle=white,toptitle=0.3mm,bottomtitle=0.3mm}
\newtcolorbox{hinweis}{enhanced,breakable,colback=yellow!12,colframe=orange!70!black,boxrule=0.6pt,arc=2mm,left=2mm,right=2mm,top=1mm,bottom=1mm,before skip=2mm,after skip=2mm}
\newcommand{\web}[1]{{\small\color{gray}\textit{Website: #1}}}

% ---------- Lücken, Linien, Karos ----------
\newcommand{\luecke}[1][3.5cm]{\rule[-0.5ex]{#1}{0.4pt}}
\newcommand{\linien}[1]{\par\vspace{1.5mm}\foreach \i in {1,...,#1}{\noindent\rule[0pt]{\linewidth}{0.3pt}\par\vspace{5.5mm}}\vspace{-3mm}}
\newcommand{\kariert}[1]{\par\vspace{1.5mm}\noindent\begin{tikzpicture}\draw[step=5mm,gray!55,thin] (0,0) grid (\linewidth-1pt,#1*5mm);\end{tikzpicture}\par}
\newcommand{\karobox}[2]{\begin{tikzpicture}\draw[step=5mm,gray!55,thin] (0,0) grid (#1,#2);\end{tikzpicture}}
\newcommand{\code}[1]{\texttt{#1}}
% Schritt-Nummern wie im Code und auf der Website: \nr{1} = ①
\newcommand{\nr}[1]{\ding{\numexpr171+#1\relax}}

% Lücke mit Lösung: \lsg[Mindestbreite]{Lösungstext}
%  Die Lösung steht (ohne eigene Höhe/Tiefe) auf derselben Linie wie im Blatt: Das Layout verschiebt sich nicht,
%  solange der Text nicht breiter als die Mindestbreite ist.
\newlength{\lsgbreite}
\newcommand{\lsg}[2][3.5cm]{\leavevmode\ifloesung\settowidth{\lsgbreite}{\,#2\,}\ifdim\lsgbreite<#1\setlength{\lsgbreite}{#1}\fi\rlap{\smash{\raisebox{0.8pt}{\makebox[\lsgbreite][l]{\lsgfarbe\,#2}}}}\luecke[\lsgbreite]\else\luecke[#1]\fi}
% Schreiblinien mit Lösung: \lsglinien{Anzahl}{Lösungstext} (gleiche Höhe wie \linien:
%  je Linie \baselineskip + 5,5 mm + \parskip, ausgemessen)
\newlength{\lsglinienhoehe}
\newcommand{\lsglinien}[2]{\ifloesung\par\vspace{1.5mm}\setlength{\lsglinienhoehe}{\dimexpr #1\dimexpr\baselineskip+5.5mm+\parskip\relax-16.25pt\relax}\noindent\begin{minipage}[t][\lsglinienhoehe][t]{\linewidth}\lsgfarbe #2\end{minipage}\par\vspace{-1.5mm}\else\linien{#1}\fi}
% Beliebiger Inhalt: \lsgoder{nur in der Lösung}{nur im Blatt}
\newcommand{\lsgoder}[2]{\ifloesung\leavevmode#1\else#2\fi}
% Text nur in der Lösung sichtbar, belegt im Blatt aber denselben Platz (ohne Text im PDF)
\newcommand{\lsgtext}[1]{\leavevmode\ifloesung{\lsgfarbe#1}\else\phantom{#1}\fi}
% Ankreuzfeld: \kreuz{1} = angekreuzt (nur in der Lösung), \kreuz{0} = leer
\newcommand{\kreuz}[1]{\leavevmode\ifloesung\ifnum#1=1 {\lsgfarbe$\boxtimes$}\else$\square$\fi\else$\square$\fi}
% Tabellenzeile mit fester Höhe für Lösungseinträge (2 Zeilen \small)
\newcommand{\zeile}{\rule[-6mm]{0pt}{9.5mm}}

% ---------- Verkettete Strukturen zeichnen ----------
% Objekt mit Kopf und Attributen:  \node[obj] (k1) {k1: Knoten \nodepart{second} daten = Hans\\ nachfolger};
\tikzset{
  obj/.style={draw,rectangle split,rectangle split parts=2,rectangle split part fill={orange!25,orange!8},
              rounded corners=2pt,font=\scriptsize,inner sep=2pt,align=left},
  qobj/.style={obj,rectangle split part fill={teal!25,teal!8}},
  aobj/.style={obj,rectangle split part fill={gray!30,gray!8}},
  pobj/.style={draw,rounded corners=2pt,fill=gray!8,font=\scriptsize,inner sep=2pt,align=left},
  klasse/.style={draw,rectangle split,rectangle split parts=3,font=\scriptsize,align=left,fill=white},
  zeiger/.style={thick,-{Stealth[length=2.2mm]}},
  erbt/.style={thick,-{Triangle[open,length=2.6mm,width=2.4mm]}},
  implementiert/.style={thick,densely dashed,-{Triangle[open,length=2.6mm,width=2.4mm]}},
  lsgzeiger/.style={zeiger,red!75!black},
  nullref/.style={font=\scriptsize\ttfamily,inner sep=1pt},
  schritt/.style={circle,draw,fill=white,inner sep=0.6pt,font=\scriptsize\bfseries},
  lsgschritt/.style={schritt,draw=red!75!black,text=red!75!black},
  % Sequenzdiagramm (UML 2): Kopf eckig, synchron = gefüllte Spitze, asynchron = offene Spitze,
  %  Rückgabe und Erzeugung gestrichelt mit offener Spitze (kein hohles Dreieck: das wäre Vererbung)
  lebenslinie/.style={draw,sharp corners,fill=white,font=\scriptsize,inner sep=2pt},
  leben/.style={densely dashed,gray},
  aufruf/.style={thick,-{Stealth[length=2mm]}},
  asynchron/.style={thick,-{Straight Barb[length=2mm,width=2.4mm]}},
  rueckgabe/.style={thick,densely dashed,-{Straight Barb[length=2mm,width=2.4mm]}},
  erzeugung/.style={rueckgabe},
}
% Akteur (Strichmännchen) mit Namen darunter: \akteur{x}{y oben}{Name}; ca. 0,9 cm hoch
\newcommand{\akteur}[3]{\begin{scope}[shift={(#1,#2)},line width=0.6pt]
  \draw (0,-0.09cm) circle (0.09cm);
  \draw (0,-0.18cm) -- (0,-0.42cm) (-0.16cm,-0.26cm) -- (0.16cm,-0.26cm)
        (-0.13cm,-0.62cm) -- (0,-0.42cm) -- (0.13cm,-0.62cm);
  \node[font=\scriptsize,inner sep=1pt,anchor=north] at (0,-0.63cm) {#3};
\end{scope}}
% Aktivierungsbalken auf der Lebenslinie bei x von y1 bis y2: \aktiv{x}{y1}{y2}
\newcommand{\aktiv}[3]{\draw[fill=gray!15] ($(#1,#2)+(-0.09cm,0)$) rectangle ($(#1,#3)+(0.09cm,0)$);}
% Lösungszeichnung: im Blatt unsichtbar, belegt aber denselben Platz
\ifloesung
  \tikzset{lsgbild/.style={red!75!black}}
\else
  \tikzset{lsgbild/.style={draw opacity=0,fill opacity=0,text opacity=0}}
\fi

% ---------- Java-Listings ----------
\lstset{language=Java,basicstyle=\ttfamily\small,keywordstyle=\bfseries,commentstyle=\color{gray!80!black},stringstyle=\color{black},showstringspaces=false,columns=flexible,keepspaces=true,tabsize=3,frame=single,framerule=0.4pt,rulecolor=\color{black},xleftmargin=2mm,framexleftmargin=1mm,aboveskip=2mm,belowskip=2mm,breaklines=true,escapechar=§,
 morekeywords={interface,implements,abstract,extends},
 literate={ä}{{\"a}}1 {ö}{{\"o}}1 {ü}{{\"u}}1 {Ä}{{\"A}}1 {Ö}{{\"O}}1 {Ü}{{\"U}}1 {ß}{{\ss}}1 {–}{{--}}1 {„}{{\glqq}}1 {“}{{\grqq}}1 {→}{{$\to$}}1}
\lstnewenvironment{java}[1][]{\lstset{#1}}{}

%  Übersetzen mit XeLaTeX, LuaLaTeX, pdfLaTeX oder Tectonic.
%
%  Blatt:   xelatex ab-1-2-interfaces.tex
%  Lösung:  xelatex ab-1-2-interfaces-lsg.tex
%           (Datei enthält nur: \def\mitloesung{}\documentclass[11pt,a4paper]{article}
% =====================================================================
%  Gemeinsame Vorlage der Arbeitsblätter "Rekursive Datenstrukturen"
%  – Informatik 12 (gleiches Layout wie informatik/13/formale-sprachen/
%  arbeitsblaetter/ab-vorlage.tex, dazu Listings und Objektdiagramme)
%  Einbinden in jedem Blatt direkt nach \documentclass: \input{ab-vorlage}
%  Übersetzen mit XeLaTeX, LuaLaTeX, pdfLaTeX oder Tectonic.
%
%  Blatt:   xelatex ab-1-2-interfaces.tex
%  Lösung:  xelatex ab-1-2-interfaces-lsg.tex
%           (Datei enthält nur: \def\mitloesung{}\input{ab-1-2-interfaces};
%            füllt alle Lücken rot aus; siehe \lsg, \lsglinien, \lsgoder,
%            \lsgtext, lsgbild, \kreuz)
% =====================================================================
\usepackage[a4paper,left=16mm,right=16mm,top=23mm,bottom=13mm,headheight=12mm,headsep=3mm,footskip=6mm]{geometry}
\usepackage{iftex}
\ifPDFTeX
  \usepackage[T1]{fontenc}
  \usepackage[utf8]{inputenc}
  \usepackage{helvet}
  \usepackage[scaled=0.86]{beramono}
  \renewcommand{\familydefault}{\sfdefault}
\else
  \usepackage{fontspec}
  \setmainfont{texgyreheros}[Extension=.otf,UprightFont=*-regular,BoldFont=*-bold,ItalicFont=*-italic,BoldItalicFont=*-bolditalic]
  \setmonofont{DejaVuSansMono}[Extension=.ttf,UprightFont=*,BoldFont=*-Bold,ItalicFont=*-Oblique,BoldItalicFont=*-BoldOblique,Scale=0.86]
\fi
\usepackage[ngerman,shorthands=off]{babel}
\usepackage{amssymb}
\usepackage{xcolor,graphicx,tabularx,array,enumitem,fancyhdr,tikz,listings,multicol,calc,etoolbox,pifont}
\usepackage[most]{tcolorbox}
\usetikzlibrary{arrows.meta,positioning,calc,decorations.pathreplacing,shapes.multipart}
\usepackage[hidelinks]{hyperref}
\setlength{\parindent}{0pt}
\setlength{\parskip}{3pt}
\setlist{nosep,leftmargin=*}

% ---------- Lösungsschalter ----------
\newif\ifloesung
\ifdefined\mitloesung\loesungtrue\fi
\newcommand{\lsgfarbe}{\color{red!75!black}}

% ---------- Kopf- und Fußzeile (einheitliche Form) ----------
\newcommand{\lehrkraft}{Hau}
\newcommand{\themenbereich}{1. Rekursive Datenstrukturen}
\newcommand{\abnummer}{}
\pagestyle{fancy}\fancyhf{}
\renewcommand{\headrulewidth}{0.4pt}
\fancyhead[L]{\small Klasse 12\\Datum:}
\fancyhead[C]{\small Themenbereich:\\\themenbereich}
\fancyhead[R]{\small \lehrkraft\ifloesung\\{\lsgfarbe\bfseries Lösung}\fi}
\fancyfoot[C]{\scriptsize Informatik 12\,·\,Blatt \abnummer\ifloesung\ (Lösung)\fi\,·\,Seite \thepage}
\newcommand{\abtitel}[2]{\renewcommand{\abnummer}{#1}\begin{center}\Large\bfseries #1\quad\underline{#2}\end{center}\vspace{-2mm}}

% ---------- Boxen ----------
\newenvironment{aufgabe}[2][]{\begin{tcolorbox}[enhanced,breakable,colback=white,colframe=black,boxrule=0.7pt,arc=3mm,left=2mm,right=2mm,top=1mm,bottom=1.2mm,before skip=2mm,after skip=2mm]\textbf{Aufgabe #2}\ifstrempty{#1}{}{ \textit{#1}}\par\vspace{0.3mm}}{\end{tcolorbox}}
\newtcolorbox{merke}[1][Merke]{enhanced,breakable,colback=green!5,colframe=green!40!black,boxrule=0.8pt,arc=2mm,left=2mm,right=2mm,top=1mm,bottom=1.2mm,before skip=2mm,after skip=2mm,title={#1},fonttitle=\bfseries,coltitle=white,toptitle=0.3mm,bottomtitle=0.3mm}
\newtcolorbox{hinweis}{enhanced,breakable,colback=yellow!12,colframe=orange!70!black,boxrule=0.6pt,arc=2mm,left=2mm,right=2mm,top=1mm,bottom=1mm,before skip=2mm,after skip=2mm}
\newcommand{\web}[1]{{\small\color{gray}\textit{Website: #1}}}

% ---------- Lücken, Linien, Karos ----------
\newcommand{\luecke}[1][3.5cm]{\rule[-0.5ex]{#1}{0.4pt}}
\newcommand{\linien}[1]{\par\vspace{1.5mm}\foreach \i in {1,...,#1}{\noindent\rule[0pt]{\linewidth}{0.3pt}\par\vspace{5.5mm}}\vspace{-3mm}}
\newcommand{\kariert}[1]{\par\vspace{1.5mm}\noindent\begin{tikzpicture}\draw[step=5mm,gray!55,thin] (0,0) grid (\linewidth-1pt,#1*5mm);\end{tikzpicture}\par}
\newcommand{\karobox}[2]{\begin{tikzpicture}\draw[step=5mm,gray!55,thin] (0,0) grid (#1,#2);\end{tikzpicture}}
\newcommand{\code}[1]{\texttt{#1}}
% Schritt-Nummern wie im Code und auf der Website: \nr{1} = ①
\newcommand{\nr}[1]{\ding{\numexpr171+#1\relax}}

% Lücke mit Lösung: \lsg[Mindestbreite]{Lösungstext}
%  Die Lösung steht (ohne eigene Höhe/Tiefe) auf derselben Linie wie im Blatt: Das Layout verschiebt sich nicht,
%  solange der Text nicht breiter als die Mindestbreite ist.
\newlength{\lsgbreite}
\newcommand{\lsg}[2][3.5cm]{\leavevmode\ifloesung\settowidth{\lsgbreite}{\,#2\,}\ifdim\lsgbreite<#1\setlength{\lsgbreite}{#1}\fi\rlap{\smash{\raisebox{0.8pt}{\makebox[\lsgbreite][l]{\lsgfarbe\,#2}}}}\luecke[\lsgbreite]\else\luecke[#1]\fi}
% Schreiblinien mit Lösung: \lsglinien{Anzahl}{Lösungstext} (gleiche Höhe wie \linien:
%  je Linie \baselineskip + 5,5 mm + \parskip, ausgemessen)
\newlength{\lsglinienhoehe}
\newcommand{\lsglinien}[2]{\ifloesung\par\vspace{1.5mm}\setlength{\lsglinienhoehe}{\dimexpr #1\dimexpr\baselineskip+5.5mm+\parskip\relax-16.25pt\relax}\noindent\begin{minipage}[t][\lsglinienhoehe][t]{\linewidth}\lsgfarbe #2\end{minipage}\par\vspace{-1.5mm}\else\linien{#1}\fi}
% Beliebiger Inhalt: \lsgoder{nur in der Lösung}{nur im Blatt}
\newcommand{\lsgoder}[2]{\ifloesung\leavevmode#1\else#2\fi}
% Text nur in der Lösung sichtbar, belegt im Blatt aber denselben Platz (ohne Text im PDF)
\newcommand{\lsgtext}[1]{\leavevmode\ifloesung{\lsgfarbe#1}\else\phantom{#1}\fi}
% Ankreuzfeld: \kreuz{1} = angekreuzt (nur in der Lösung), \kreuz{0} = leer
\newcommand{\kreuz}[1]{\leavevmode\ifloesung\ifnum#1=1 {\lsgfarbe$\boxtimes$}\else$\square$\fi\else$\square$\fi}
% Tabellenzeile mit fester Höhe für Lösungseinträge (2 Zeilen \small)
\newcommand{\zeile}{\rule[-6mm]{0pt}{9.5mm}}

% ---------- Verkettete Strukturen zeichnen ----------
% Objekt mit Kopf und Attributen:  \node[obj] (k1) {k1: Knoten \nodepart{second} daten = Hans\\ nachfolger};
\tikzset{
  obj/.style={draw,rectangle split,rectangle split parts=2,rectangle split part fill={orange!25,orange!8},
              rounded corners=2pt,font=\scriptsize,inner sep=2pt,align=left},
  qobj/.style={obj,rectangle split part fill={teal!25,teal!8}},
  aobj/.style={obj,rectangle split part fill={gray!30,gray!8}},
  pobj/.style={draw,rounded corners=2pt,fill=gray!8,font=\scriptsize,inner sep=2pt,align=left},
  klasse/.style={draw,rectangle split,rectangle split parts=3,font=\scriptsize,align=left,fill=white},
  zeiger/.style={thick,-{Stealth[length=2.2mm]}},
  erbt/.style={thick,-{Triangle[open,length=2.6mm,width=2.4mm]}},
  implementiert/.style={thick,densely dashed,-{Triangle[open,length=2.6mm,width=2.4mm]}},
  lsgzeiger/.style={zeiger,red!75!black},
  nullref/.style={font=\scriptsize\ttfamily,inner sep=1pt},
  schritt/.style={circle,draw,fill=white,inner sep=0.6pt,font=\scriptsize\bfseries},
  lsgschritt/.style={schritt,draw=red!75!black,text=red!75!black},
  % Sequenzdiagramm (UML 2): Kopf eckig, synchron = gefüllte Spitze, asynchron = offene Spitze,
  %  Rückgabe und Erzeugung gestrichelt mit offener Spitze (kein hohles Dreieck: das wäre Vererbung)
  lebenslinie/.style={draw,sharp corners,fill=white,font=\scriptsize,inner sep=2pt},
  leben/.style={densely dashed,gray},
  aufruf/.style={thick,-{Stealth[length=2mm]}},
  asynchron/.style={thick,-{Straight Barb[length=2mm,width=2.4mm]}},
  rueckgabe/.style={thick,densely dashed,-{Straight Barb[length=2mm,width=2.4mm]}},
  erzeugung/.style={rueckgabe},
}
% Akteur (Strichmännchen) mit Namen darunter: \akteur{x}{y oben}{Name}; ca. 0,9 cm hoch
\newcommand{\akteur}[3]{\begin{scope}[shift={(#1,#2)},line width=0.6pt]
  \draw (0,-0.09cm) circle (0.09cm);
  \draw (0,-0.18cm) -- (0,-0.42cm) (-0.16cm,-0.26cm) -- (0.16cm,-0.26cm)
        (-0.13cm,-0.62cm) -- (0,-0.42cm) -- (0.13cm,-0.62cm);
  \node[font=\scriptsize,inner sep=1pt,anchor=north] at (0,-0.63cm) {#3};
\end{scope}}
% Aktivierungsbalken auf der Lebenslinie bei x von y1 bis y2: \aktiv{x}{y1}{y2}
\newcommand{\aktiv}[3]{\draw[fill=gray!15] ($(#1,#2)+(-0.09cm,0)$) rectangle ($(#1,#3)+(0.09cm,0)$);}
% Lösungszeichnung: im Blatt unsichtbar, belegt aber denselben Platz
\ifloesung
  \tikzset{lsgbild/.style={red!75!black}}
\else
  \tikzset{lsgbild/.style={draw opacity=0,fill opacity=0,text opacity=0}}
\fi

% ---------- Java-Listings ----------
\lstset{language=Java,basicstyle=\ttfamily\small,keywordstyle=\bfseries,commentstyle=\color{gray!80!black},stringstyle=\color{black},showstringspaces=false,columns=flexible,keepspaces=true,tabsize=3,frame=single,framerule=0.4pt,rulecolor=\color{black},xleftmargin=2mm,framexleftmargin=1mm,aboveskip=2mm,belowskip=2mm,breaklines=true,escapechar=§,
 morekeywords={interface,implements,abstract,extends},
 literate={ä}{{\"a}}1 {ö}{{\"o}}1 {ü}{{\"u}}1 {Ä}{{\"A}}1 {Ö}{{\"O}}1 {Ü}{{\"U}}1 {ß}{{\ss}}1 {–}{{--}}1 {„}{{\glqq}}1 {“}{{\grqq}}1 {→}{{$\to$}}1}
\lstnewenvironment{java}[1][]{\lstset{#1}}{}

\begin{document}
\abtitel{1.2}{Interfaces: ein Vertrag für alle Daten}

Am Bahnhof warten Taxis in einer Schlange: Wer zuerst ankommt, fährt zuerst. Die Stadt möchte dafür die \code{Warteschlange} aus 1.1 benutzen. Doch die kann bisher nur Patienten verwalten. Eine Kopie für Taxis (\code{TaxiKnoten}, \code{TaxiWarteschlange}) würde funktionieren, aber jede Verbesserung müsste man dann doppelt einbauen. \web{Station 1.2, Projekt Taxistand}

\begin{aufgabe}[Was muss die Warteschlange über ihre Daten wissen?]{1}
Untersuche im Projekt die Methoden der Warteschlange. Kreuze an, ob die Methode die Daten nur weiterreicht oder eine Methode der Daten aufruft. Notiere diese gegebenenfalls.\\[1mm]
\small\renewcommand{\arraystretch}{1.2}
\begin{tabularx}{\linewidth}{|l|c|c|X|}\hline
\textbf{Methode} & \textbf{reicht nur weiter} & \textbf{ruft eine Methode der Daten auf} & \textbf{welche?}\\\hline
\code{add(Patient p)} & \kreuz{1} & \kreuz{0} & \\\hline
\code{poll()}, \code{peek()} & \kreuz{1} & \kreuz{0} & \\\hline
\code{ausgeben()} & \kreuz{0} & \kreuz{1} & \lsgtext{\code{getName()}: der Name zum Ausgeben}\\\hline
\end{tabularx}\\[1.5mm]
\normalsize
\textbf{Folgerung:} Von ihren Daten muss die Warteschlange nur verlangen, dass sie einen \lsg[4.6cm]{Namen (Schlüsselwert) nennen} können. Genau das legen wir in einem Vertrag fest: im \textbf{Interface} \code{DatenElement}.
\end{aufgabe}

\begin{minipage}[t]{0.47\linewidth}
\vspace{0pt}
\begin{java}[basicstyle=\ttfamily\footnotesize]
public §\lsg[1.9cm]{interface}§ DatenElement {
   String gibSchluesselwert();
}

public class Taxi §\lsg[3.2cm]{implements DatenElement}§ {
   private int nummer;
   private String fahrer;
   ...
   public String gibSchluesselwert() {
      return §\lsg[2.6cm]{"Taxi " + nummer}§;
   }
}
\end{java}
\end{minipage}\hfill
\begin{minipage}[t]{0.505\linewidth}
\vspace{0pt}
\begin{merke}[Interface]
\small
\begin{itemize}
\item Ein Interface ist ein \lsg[1.7cm]{Vertrag}. Es legt nur \lsg[2.9cm]{Methodenköpfe} fest, keine Rümpfe.
\item Eine Klasse, die es implementiert (\code{implements}), muss \lsg[1.3cm]{alle} seine Methoden bereitstellen. Das prüft der \lsg[2cm]{Compiler}.
\item Ein Interface ist auch ein \lsg[1.9cm]{Datentyp}: Eine Variable vom Typ \code{DatenElement} kann Objekte jeder Klasse aufnehmen, die es implementiert. Objekte vom Interface selbst gibt es \lsg[1.2cm]{nicht}.
\end{itemize}
\end{merke}
\end{minipage}

\begin{aufgabe}[Klassendiagramm]{2}
\textbf{a)} Zeichne alle Beziehungen ein: Assoziationen mit Namen und Multiplizitäten, dazu \code{implements} als gestrichelten Pfeil mit leerem Dreieck zum Interface.\\
\textbf{b)} Umrande die Klassen der \textbf{Struktur} blau und die Klassen des \textbf{Inhalts} grün.\\[1mm]
\begin{tikzpicture}
\node[klasse,fill=teal!6,anchor=north west] (W) at (0,0)
  {\textbf{Warteschlange}\nodepart{second}- anfang: Knoten\\- ende: Knoten
   \nodepart{third}+ add(DatenElement d): boolean\\+ poll(): DatenElement\\+ ausgeben()};
\node[klasse,fill=orange!8,anchor=north west] (K) at (6.0,0)
  {\textbf{Knoten}\nodepart{second}- daten: DatenElement\\- nachfolger: Knoten
   \nodepart{third}+ getDaten(): DatenElement\\+ getNachfolger(): Knoten};
\node[klasse,fill=green!6,anchor=north,rectangle split parts=2] (D) at (13.1,0)
  {\textit{\guillemotleft interface\guillemotright}\\\textbf{DatenElement}\nodepart{second}+ gibSchluesselwert(): String};
\node[klasse,fill=gray!8,anchor=north] (P) at (11.2,-3.0)
  {\textbf{Patient}\nodepart{second}- name: String
   \nodepart{third}+ gibSchluesselwert(): String};
\node[klasse,fill=gray!8,anchor=north] (T) at (15.15,-3.0)
  {\textbf{Taxi}\nodepart{second}- nummer: int\\- fahrer: String
   \nodepart{third}+ gibSchluesselwert(): String};
\begin{scope}[font=\scriptsize]
  \draw[zeiger,lsgbild] ([yshift=-4mm]W.north east) -- ([yshift=-4mm]W.north east -| K.west)
     node[pos=0.3,above] {\lsgtext{anfang}} node[pos=0.86,below] {\lsgtext{0..1}};
  \draw[zeiger,lsgbild] ([yshift=-13mm]W.north east) -- ([yshift=-13mm]W.north east -| K.west)
     node[pos=0.3,below] {\lsgtext{ende}} node[pos=0.86,above] {\lsgtext{0..1}};
  \draw[zeiger,lsgbild] ([yshift=-5mm]K.north east) -- ([yshift=-5mm]K.north east -| D.west)
     node[pos=0.6,below] {\lsgtext{daten}} node[pos=0.88,above] {\lsgtext{1}};
  \draw[zeiger,lsgbild] ([xshift=-5mm]K.north east) -- ++(0,4mm) -| ([xshift=5mm]K.north west);
  \node[left] at ($(K.north west)+(4.5mm,3.2mm)$) {\lsgtext{nachfolger\quad 0..1}};
  \draw[implementiert,lsgbild] ([xshift=10mm]P.north) -- ([xshift=10mm]P.north |- D.south);
  \draw[implementiert,lsgbild] ([xshift=-9.5mm]T.north) -- ([xshift=-9.5mm]T.north |- D.south);
  \draw[lsgbild,blue!70!black,thick,rounded corners] ($(W.north west)+(-1.5mm,7.5mm)$) rectangle ($(K.east |- W.south)+(1.5mm,-1.5mm)$);
  % Inhalt: oben nur um das Interface, unten um Patient und Taxi (unter dem blauen Rahmen)
  \coordinate (gl) at ($(K.east)+(2.75mm,0)$);
  \coordinate (gt) at ($(D.north)+(0,2mm)$);
  \coordinate (gr) at ($(T.east)+(1.5mm,0)$);
  \coordinate (gb) at ($(T.south)+(0,-1.5mm)$);
  \coordinate (gp) at ($(P.west)+(-1.5mm,0)$);
  \coordinate (gm) at ($(W.south)+(0,-3mm)$);
  \draw[lsgbild,green!45!black,thick,rounded corners] (gl |- gt) -- (gr |- gt) -- (gr |- gb) -- (gp |- gb) -- (gp |- gm) -- (gl |- gm) -- cycle;
\end{scope}
\end{tikzpicture}\\[1mm]
\textbf{c)} Durch die Trennung von Struktur und Inhalt erledigt jede Klasse genau eine Aufgabe (hohe \lsg[2.1cm]{Kohäsion}). Struktur und Inhalt hängen nur über das Interface zusammen (lose \lsg[2.1cm]{Kopplung}). So kann die Warteschlange ohne jede Änderung \lsg[4.5cm]{jedes neue DatenElement} aufnehmen.
\end{aufgabe}

\newpage
\begin{aufgabe}[Projekt Taxistand, TODO 3a bis 3d]{3}
Baue das Projekt so um, dass Warteschlange und Knoten mit jedem \code{DatenElement} arbeiten. Ergänze die Tabelle und übertrage die Änderungen ins Projekt (Reihenfolge wie die TODOs).\\[1mm]
\small\renewcommand{\arraystretch}{1.25}
\begin{tabularx}{\linewidth}{|c|l|X|}\hline
 & \textbf{bisher} & \textbf{neu}\\\hline
3a & \code{DatenElement}: (leer) & \lsgtext{\code{String gibSchluesselwert();}}\\\hline
3b & \code{public class Patient \{} & \lsgtext{\code{public class Patient implements DatenElement \{}}\\
   & & \lsgtext{\code{+ public String gibSchluesselwert() \{ return name; \}}}\\\hline
3c & \code{private Patient daten;} & \lsgtext{\code{private DatenElement daten;}} \ (ebenso Konstruktor, \code{getDaten})\\\hline
3d & \code{public boolean add(Patient p)} & \lsgtext{\code{public boolean add(DatenElement daten)}}\\
   & \code{public Patient poll()} & \lsgtext{\code{public DatenElement poll()}} \ (ebenso \code{peek})\\
   & \code{aktuell.getDaten().getName()} & \lsgtext{\code{aktuell.getDaten().gibSchluesselwert()}}\\\hline
\end{tabularx}\\[1.5mm]
\normalsize
\textbf{b)} Nach dem Umbau von \code{Knoten} meldet der Compiler einen Fehler in \code{ausgeben()}. Erkläre!\\
\lsg[\linewidth]{Ein DatenElement hat nur die Methode gibSchluesselwert(), kein getName().}\\[1mm]
\lsg[\linewidth]{Die Warteschlange darf nur benutzen, was im Vertrag steht.}
\end{aufgabe}

\begin{aufgabe}[Projekt Taxistand, TODO 4]{4}
\textbf{a)} Ergänze den Code von \code{Taxi} oben auf Seite 1, übertrage ihn ins Projekt und kommentiere im Hauptprogramm die Taxi-Zeilen ein.\\
\textbf{b)} Lösche testweise die Methode \code{gibSchluesselwert()} aus \code{Taxi}. Notiere die Meldung des Compilers und erkläre sie mit dem Wort \textit{Vertrag}.\\
\lsg[\linewidth]{„Die Klasse Taxi muss noch folgende Methoden des Interfaces DatenElement implementieren: …“}\\[1mm]
\lsg[\linewidth]{Taxi verspricht mit implements, den Vertrag zu erfüllen, hält ihn aber nicht ein.}
\end{aufgabe}

\begin{aufgabe}[Datentypen]{5}
\code{taxistand} ist eine \code{Warteschlange}. Kreuze an, welche Zeilen der Compiler akzeptiert, und begründe kurz.\\[1mm]
\small\renewcommand{\arraystretch}{1.25}
\begin{tabularx}{\linewidth}{|l|c|c|X|}\hline
\textbf{Zeile} & \textbf{ok} & \textbf{Fehler} & \textbf{Begründung}\\\hline
\code{DatenElement d = new Taxi(5, "Lea");} & \kreuz{1} & \kreuz{0} & \lsgtext{Taxi implementiert DatenElement.}\\\hline
\code{DatenElement d = new DatenElement();} & \kreuz{0} & \kreuz{1} & \lsgtext{Von einem Interface gibt es keine Objekte.}\\\hline
\code{taxistand.add(new Patient("Ali", "Husten"));} & \kreuz{1} & \kreuz{0} & \lsgtext{Auch Patient ist ein DatenElement.}\\\hline
\code{Taxi t = taxistand.poll();} & \kreuz{0} & \kreuz{1} & \lsgtext{poll() liefert nur ein DatenElement.}\\\hline
\code{println(taxistand.poll().gibSchluesselwert());} & \kreuz{1} & \kreuz{0} & \lsgtext{Diese Methode steht im Vertrag.}\\\hline
\code{println(taxistand.poll().getFahrer());} & \kreuz{0} & \kreuz{1} & \lsgtext{getFahrer() steht nicht im Vertrag.}\\\hline
\end{tabularx}
\normalsize
\end{aufgabe}

\begin{hinweis}\small
\textbf{Es geht auch anders -- besprecht mündlich, warum man es so machen kann:}
\begin{itemize}
\item Jede Klasse ist ein \code{Object}. Der Knoten könnte also einfach \code{Object daten} speichern. Was ginge dann in \code{ausgeben()} nicht mehr?
\item In einer älteren Fassung stand im Vertrag nur \code{String toString();}. Warum erzwingt dieser Vertrag nichts? (Tipp: Was erbt jede Klasse von \code{Object}?)
\item Statt des Interfaces gäbe es eine Oberklasse \code{Daten}, von der Patient und Taxi erben. Was spricht dagegen, wenn \code{Patient} schon von \code{Person} erbt?
\item Du kennst schon ein Interface: \code{Wartezimmer} aus 1.1. Was legte es fest -- und wozu?
\end{itemize}
\end{hinweis}

\end{document}
;
%            füllt alle Lücken rot aus; siehe \lsg, \lsglinien, \lsgoder,
%            \lsgtext, lsgbild, \kreuz)
% =====================================================================
\usepackage[a4paper,left=16mm,right=16mm,top=23mm,bottom=13mm,headheight=12mm,headsep=3mm,footskip=6mm]{geometry}
\usepackage{iftex}
\ifPDFTeX
  \usepackage[T1]{fontenc}
  \usepackage[utf8]{inputenc}
  \usepackage{helvet}
  \usepackage[scaled=0.86]{beramono}
  \renewcommand{\familydefault}{\sfdefault}
\else
  \usepackage{fontspec}
  \setmainfont{texgyreheros}[Extension=.otf,UprightFont=*-regular,BoldFont=*-bold,ItalicFont=*-italic,BoldItalicFont=*-bolditalic]
  \setmonofont{DejaVuSansMono}[Extension=.ttf,UprightFont=*,BoldFont=*-Bold,ItalicFont=*-Oblique,BoldItalicFont=*-BoldOblique,Scale=0.86]
\fi
\usepackage[ngerman,shorthands=off]{babel}
\usepackage{amssymb}
\usepackage{xcolor,graphicx,tabularx,array,enumitem,fancyhdr,tikz,listings,multicol,calc,etoolbox,pifont}
\usepackage[most]{tcolorbox}
\usetikzlibrary{arrows.meta,positioning,calc,decorations.pathreplacing,shapes.multipart}
\usepackage[hidelinks]{hyperref}
\setlength{\parindent}{0pt}
\setlength{\parskip}{3pt}
\setlist{nosep,leftmargin=*}

% ---------- Lösungsschalter ----------
\newif\ifloesung
\ifdefined\mitloesung\loesungtrue\fi
\newcommand{\lsgfarbe}{\color{red!75!black}}

% ---------- Kopf- und Fußzeile (einheitliche Form) ----------
\newcommand{\lehrkraft}{Hau}
\newcommand{\themenbereich}{1. Rekursive Datenstrukturen}
\newcommand{\abnummer}{}
\pagestyle{fancy}\fancyhf{}
\renewcommand{\headrulewidth}{0.4pt}
\fancyhead[L]{\small Klasse 12\\Datum:}
\fancyhead[C]{\small Themenbereich:\\\themenbereich}
\fancyhead[R]{\small \lehrkraft\ifloesung\\{\lsgfarbe\bfseries Lösung}\fi}
\fancyfoot[C]{\scriptsize Informatik 12\,·\,Blatt \abnummer\ifloesung\ (Lösung)\fi\,·\,Seite \thepage}
\newcommand{\abtitel}[2]{\renewcommand{\abnummer}{#1}\begin{center}\Large\bfseries #1\quad\underline{#2}\end{center}\vspace{-2mm}}

% ---------- Boxen ----------
\newenvironment{aufgabe}[2][]{\begin{tcolorbox}[enhanced,breakable,colback=white,colframe=black,boxrule=0.7pt,arc=3mm,left=2mm,right=2mm,top=1mm,bottom=1.2mm,before skip=2mm,after skip=2mm]\textbf{Aufgabe #2}\ifstrempty{#1}{}{ \textit{#1}}\par\vspace{0.3mm}}{\end{tcolorbox}}
\newtcolorbox{merke}[1][Merke]{enhanced,breakable,colback=green!5,colframe=green!40!black,boxrule=0.8pt,arc=2mm,left=2mm,right=2mm,top=1mm,bottom=1.2mm,before skip=2mm,after skip=2mm,title={#1},fonttitle=\bfseries,coltitle=white,toptitle=0.3mm,bottomtitle=0.3mm}
\newtcolorbox{hinweis}{enhanced,breakable,colback=yellow!12,colframe=orange!70!black,boxrule=0.6pt,arc=2mm,left=2mm,right=2mm,top=1mm,bottom=1mm,before skip=2mm,after skip=2mm}
\newcommand{\web}[1]{{\small\color{gray}\textit{Website: #1}}}

% ---------- Lücken, Linien, Karos ----------
\newcommand{\luecke}[1][3.5cm]{\rule[-0.5ex]{#1}{0.4pt}}
\newcommand{\linien}[1]{\par\vspace{1.5mm}\foreach \i in {1,...,#1}{\noindent\rule[0pt]{\linewidth}{0.3pt}\par\vspace{5.5mm}}\vspace{-3mm}}
\newcommand{\kariert}[1]{\par\vspace{1.5mm}\noindent\begin{tikzpicture}\draw[step=5mm,gray!55,thin] (0,0) grid (\linewidth-1pt,#1*5mm);\end{tikzpicture}\par}
\newcommand{\karobox}[2]{\begin{tikzpicture}\draw[step=5mm,gray!55,thin] (0,0) grid (#1,#2);\end{tikzpicture}}
\newcommand{\code}[1]{\texttt{#1}}
% Schritt-Nummern wie im Code und auf der Website: \nr{1} = ①
\newcommand{\nr}[1]{\ding{\numexpr171+#1\relax}}

% Lücke mit Lösung: \lsg[Mindestbreite]{Lösungstext}
%  Die Lösung steht (ohne eigene Höhe/Tiefe) auf derselben Linie wie im Blatt: Das Layout verschiebt sich nicht,
%  solange der Text nicht breiter als die Mindestbreite ist.
\newlength{\lsgbreite}
\newcommand{\lsg}[2][3.5cm]{\leavevmode\ifloesung\settowidth{\lsgbreite}{\,#2\,}\ifdim\lsgbreite<#1\setlength{\lsgbreite}{#1}\fi\rlap{\smash{\raisebox{0.8pt}{\makebox[\lsgbreite][l]{\lsgfarbe\,#2}}}}\luecke[\lsgbreite]\else\luecke[#1]\fi}
% Schreiblinien mit Lösung: \lsglinien{Anzahl}{Lösungstext} (gleiche Höhe wie \linien:
%  je Linie \baselineskip + 5,5 mm + \parskip, ausgemessen)
\newlength{\lsglinienhoehe}
\newcommand{\lsglinien}[2]{\ifloesung\par\vspace{1.5mm}\setlength{\lsglinienhoehe}{\dimexpr #1\dimexpr\baselineskip+5.5mm+\parskip\relax-16.25pt\relax}\noindent\begin{minipage}[t][\lsglinienhoehe][t]{\linewidth}\lsgfarbe #2\end{minipage}\par\vspace{-1.5mm}\else\linien{#1}\fi}
% Beliebiger Inhalt: \lsgoder{nur in der Lösung}{nur im Blatt}
\newcommand{\lsgoder}[2]{\ifloesung\leavevmode#1\else#2\fi}
% Text nur in der Lösung sichtbar, belegt im Blatt aber denselben Platz (ohne Text im PDF)
\newcommand{\lsgtext}[1]{\leavevmode\ifloesung{\lsgfarbe#1}\else\phantom{#1}\fi}
% Ankreuzfeld: \kreuz{1} = angekreuzt (nur in der Lösung), \kreuz{0} = leer
\newcommand{\kreuz}[1]{\leavevmode\ifloesung\ifnum#1=1 {\lsgfarbe$\boxtimes$}\else$\square$\fi\else$\square$\fi}
% Tabellenzeile mit fester Höhe für Lösungseinträge (2 Zeilen \small)
\newcommand{\zeile}{\rule[-6mm]{0pt}{9.5mm}}

% ---------- Verkettete Strukturen zeichnen ----------
% Objekt mit Kopf und Attributen:  \node[obj] (k1) {k1: Knoten \nodepart{second} daten = Hans\\ nachfolger};
\tikzset{
  obj/.style={draw,rectangle split,rectangle split parts=2,rectangle split part fill={orange!25,orange!8},
              rounded corners=2pt,font=\scriptsize,inner sep=2pt,align=left},
  qobj/.style={obj,rectangle split part fill={teal!25,teal!8}},
  aobj/.style={obj,rectangle split part fill={gray!30,gray!8}},
  pobj/.style={draw,rounded corners=2pt,fill=gray!8,font=\scriptsize,inner sep=2pt,align=left},
  klasse/.style={draw,rectangle split,rectangle split parts=3,font=\scriptsize,align=left,fill=white},
  zeiger/.style={thick,-{Stealth[length=2.2mm]}},
  erbt/.style={thick,-{Triangle[open,length=2.6mm,width=2.4mm]}},
  implementiert/.style={thick,densely dashed,-{Triangle[open,length=2.6mm,width=2.4mm]}},
  lsgzeiger/.style={zeiger,red!75!black},
  nullref/.style={font=\scriptsize\ttfamily,inner sep=1pt},
  schritt/.style={circle,draw,fill=white,inner sep=0.6pt,font=\scriptsize\bfseries},
  lsgschritt/.style={schritt,draw=red!75!black,text=red!75!black},
  % Sequenzdiagramm (UML 2): Kopf eckig, synchron = gefüllte Spitze, asynchron = offene Spitze,
  %  Rückgabe und Erzeugung gestrichelt mit offener Spitze (kein hohles Dreieck: das wäre Vererbung)
  lebenslinie/.style={draw,sharp corners,fill=white,font=\scriptsize,inner sep=2pt},
  leben/.style={densely dashed,gray},
  aufruf/.style={thick,-{Stealth[length=2mm]}},
  asynchron/.style={thick,-{Straight Barb[length=2mm,width=2.4mm]}},
  rueckgabe/.style={thick,densely dashed,-{Straight Barb[length=2mm,width=2.4mm]}},
  erzeugung/.style={rueckgabe},
}
% Akteur (Strichmännchen) mit Namen darunter: \akteur{x}{y oben}{Name}; ca. 0,9 cm hoch
\newcommand{\akteur}[3]{\begin{scope}[shift={(#1,#2)},line width=0.6pt]
  \draw (0,-0.09cm) circle (0.09cm);
  \draw (0,-0.18cm) -- (0,-0.42cm) (-0.16cm,-0.26cm) -- (0.16cm,-0.26cm)
        (-0.13cm,-0.62cm) -- (0,-0.42cm) -- (0.13cm,-0.62cm);
  \node[font=\scriptsize,inner sep=1pt,anchor=north] at (0,-0.63cm) {#3};
\end{scope}}
% Aktivierungsbalken auf der Lebenslinie bei x von y1 bis y2: \aktiv{x}{y1}{y2}
\newcommand{\aktiv}[3]{\draw[fill=gray!15] ($(#1,#2)+(-0.09cm,0)$) rectangle ($(#1,#3)+(0.09cm,0)$);}
% Lösungszeichnung: im Blatt unsichtbar, belegt aber denselben Platz
\ifloesung
  \tikzset{lsgbild/.style={red!75!black}}
\else
  \tikzset{lsgbild/.style={draw opacity=0,fill opacity=0,text opacity=0}}
\fi

% ---------- Java-Listings ----------
\lstset{language=Java,basicstyle=\ttfamily\small,keywordstyle=\bfseries,commentstyle=\color{gray!80!black},stringstyle=\color{black},showstringspaces=false,columns=flexible,keepspaces=true,tabsize=3,frame=single,framerule=0.4pt,rulecolor=\color{black},xleftmargin=2mm,framexleftmargin=1mm,aboveskip=2mm,belowskip=2mm,breaklines=true,escapechar=§,
 morekeywords={interface,implements,abstract,extends},
 literate={ä}{{\"a}}1 {ö}{{\"o}}1 {ü}{{\"u}}1 {Ä}{{\"A}}1 {Ö}{{\"O}}1 {Ü}{{\"U}}1 {ß}{{\ss}}1 {–}{{--}}1 {„}{{\glqq}}1 {“}{{\grqq}}1 {→}{{$\to$}}1}
\lstnewenvironment{java}[1][]{\lstset{#1}}{}

\begin{document}
\abtitel{1.3}{Listen und Rekursion}

Taxis melden sich in beliebiger Reihenfolge an und ab. Unsere \textbf{Liste} kennt nur ihren ersten Knoten (\code{anfang}), den Rest erledigen die Knoten. \web{Station 1.3, Projekt Taxizentrale}

\begin{aufgabe}[Wie viele stehen hinter dir?]{1}
\begin{minipage}[t]{0.56\linewidth}
\vspace{0pt}\raggedright
Die erste Person in der Mensa-Schlange will wissen, wie viele anstehen. Jede darf nur ihren Hintermann fragen. \textbf{a)} Ergänze die Regel:\\[1mm]
„Steht hinter mir \lsg[2.1cm]{niemand}, antworte ich \lsg[0.9cm]{1}. Sonst frage ich \lsg[3.2cm]{meinen Hintermann} und antworte mit \lsg[3.9cm]{seiner Antwort + 1}.“
\end{minipage}\hfill
\begin{minipage}[t]{0.41\linewidth}
\vspace{0pt}\raggedright
\textbf{b)} Spielt es zu fünft durch. Wer antwortet zuerst, wer zuletzt?\\
\lsg[0.98\linewidth]{zuerst die letzte Person (1),}\\[1mm]
\lsg[0.98\linewidth]{zuletzt die erste Person (5)}
\end{minipage}
\end{aufgabe}

\begin{aufgabe}[Projekt Taxizentrale, TODO 2]{2}
\textbf{a)} Ergänze \code{laenge()} nach der Regel aus Aufgabe 1. Beschrifte: \textit{Abbruchbedingung}, \textit{rekursiver Aufruf}.\\[0.5mm]
\begin{minipage}[t]{0.38\linewidth}
\vspace{0pt}
\textit{\small iterativ (Warteschlange, 1.1)}
\begin{java}[basicstyle=\ttfamily\scriptsize]
public int anzahl() {
   int zaehler = 0;
   Knoten aktuell = anfang;
   while (aktuell != null) {
      zaehler++;
      aktuell = aktuell.getNachfolger();
   }
   return zaehler;
}
\end{java}
\end{minipage}\hfill
\begin{minipage}[t]{0.595\linewidth}
\vspace{0pt}
\textit{\small rekursiv (Knoten)}
\begin{java}[basicstyle=\ttfamily\scriptsize]
public int laenge() {
   if (§\lsg[2.6cm]{nachfolger == null}§) {    // §\lsg[2.6cm]{Abbruchbedingung}§
      return §\lsg[0.6cm]{1}§;
   } else {
      return §\lsg[3.4cm]{nachfolger.laenge() + 1}§;  // §\lsg[2.6cm]{rekursiver Aufruf}§
   }
}
\end{java}
\textit{\small rekursiv (Liste)}
\begin{java}[basicstyle=\ttfamily\scriptsize]
public int laenge() {
   if (anfang == null) { return §\lsg[0.6cm]{0}§; }
   else { return §\lsg[2.6cm]{anfang.laenge()}§; }
}
\end{java}
\end{minipage}
\begin{merke}[Das UML-Sequenzdiagramm: Wer ruft wann wen auf?]
\begin{minipage}[t]{0.43\linewidth}
\vspace{0pt}
\begin{tikzpicture}[x=1cm,y=0.8cm]
\akteur{0.45}{0.25}{Nutzer:in}
\draw[leben] (0.45,-0.95) -- (0.45,-3.6);
\node[lebenslinie] (li) at (2.25,-0.55) {liste : Liste};
\node[lebenslinie] (ta) at (6.35,-0.55) {t : Taxi};
\node[lebenslinie] (kn) at (4.75,-1.85) {k1 : Knoten};
\foreach \k/\x in {li/2.25, ta/6.35, kn/4.75} { \draw[leben] (\k.south) -- (\x,-3.6); }
\aktiv{2.25}{-1.3}{-2.5}
\aktiv{6.35}{-3.05}{-3.45}
\draw[aufruf] (0.45,-1.3) -- node[above,pos=0.62,font=\scriptsize,inner sep=1pt] {\code{add(t)}} ([xshift=-0.09cm]2.25,-1.3);
\draw[erzeugung] ([xshift=0.09cm]2.25,-1.85) -- node[above,font=\scriptsize,inner sep=1pt] {«create»} (kn.west);
\draw[rueckgabe] ([xshift=-0.09cm]2.25,-2.5) -- node[above,pos=0.4,font=\scriptsize,inner sep=1pt] {\code{true}} (0.45,-2.5);
\draw[asynchron] (0.45,-3.05) -- node[above,font=\scriptsize,inner sep=1pt] {\code{funk("Bahnhof")}} ([xshift=-0.09cm]6.35,-3.05);
\foreach \x/\y/\n in {1.0/-0.3/1, 3.3/-0.55/2, 2.6/-2.17/3, 1.3/-1.65/4, 0.8/-2.24/5, 3.4/-2.22/6, 3.6/-3.4/7} { \node[schritt] at (\x,\y) {\n}; }
\end{tikzpicture}
\end{minipage}\hfill
\begin{minipage}[t]{0.55\linewidth}
\vspace{0pt}\scriptsize\raggedright
\nr{1} \textbf{Akteur}: startet den Ablauf (im Projekt: Hauptprogramm).\\
\nr{2} \textbf{Objekt} \code{name : Klasse} mit \textbf{Lebenslinie} (Zeit läuft nach unten).\\
\nr{3} \textbf{Aktivierungsbalken}: Objekt arbeitet oder wartet auf Antwort.\\
\nr{4} \textbf{Synchroner Aufruf}, gefüllte Spitze: Der Aufrufer wartet auf die Antwort (so jeder Methodenaufruf in Java).\\
\nr{5} \textbf{Rückgabe}, gestrichelt, mit Wert: Danach endet der Balken.\\
\nr{6} \textbf{Erzeugung} «create» (\code{new}): Pfeil zum Kopf des neuen Objekts.\\
\nr{7} \textbf{Asynchroner Aufruf}, offene Spitze: Der Aufrufer schickt ab und arbeitet sofort weiter, ohne Rückgabe (wie Textnachricht statt Anruf).
\end{minipage}
\end{merke}
\textbf{b)} Ergänze für \code{liste.laenge()} alle Aufrufe, Rückgaben mit Wert und Aktivierungsbalken.\\[0.5mm]
\begin{tikzpicture}[x=1cm,y=0.85cm]
\akteur{0.6}{0.5}{Nutzer:in}
\draw[leben] (0.6,-0.6) -- (0.6,-3.55);
\foreach \name/\x/\txt in {li/3.8/liste : Liste, k1/7.2/k1 : Knoten, k2/10.6/k2 : Knoten, k3/14/k3 : Knoten} {
  \node[lebenslinie] (\name) at (\x,0) {\txt};
  \draw[leben] (\name.south) -- (\x,-3.55);
}
\node[font=\tiny,gray,right=0pt of k1.south east,anchor=north west] {Taxi 12};
\node[font=\tiny,gray,right=0pt of k2.south east,anchor=north west] {Taxi 7};
\node[font=\tiny,gray,right=0pt of k3.south east,anchor=north west] {Taxi 3};
\aktiv{3.8}{-0.65}{-3.35}
\draw[aufruf] (0.6,-0.65) -- node[above,font=\scriptsize,inner sep=1pt] {\code{laenge()}} ([xshift=-0.09cm]3.8,-0.65);
\begin{scope}[lsgbild]
  \aktiv{7.2}{-1.05}{-3.0}
  \aktiv{10.6}{-1.45}{-2.65}
  \aktiv{14}{-1.85}{-2.3}
  \draw[aufruf] ([xshift=0.09cm]3.8,-1.05) -- node[above,font=\scriptsize,inner sep=1pt] {\lsgtext{\code{laenge()}}} ([xshift=-0.09cm]7.2,-1.05);
  \draw[aufruf] ([xshift=0.09cm]7.2,-1.45) -- node[above,font=\scriptsize,inner sep=1pt] {\lsgtext{\code{laenge()}}} ([xshift=-0.09cm]10.6,-1.45);
  \draw[aufruf] ([xshift=0.09cm]10.6,-1.85) -- node[above,font=\scriptsize,inner sep=1pt] {\lsgtext{\code{laenge()}}} ([xshift=-0.09cm]14,-1.85);
  \draw[rueckgabe] ([xshift=-0.09cm]14,-2.3) -- node[above,font=\scriptsize,inner sep=1pt] {\lsgtext{\code{1}}} ([xshift=0.09cm]10.6,-2.3);
  \draw[rueckgabe] ([xshift=-0.09cm]10.6,-2.65) -- node[above,font=\scriptsize,inner sep=1pt] {\lsgtext{\code{2}}} ([xshift=0.09cm]7.2,-2.65);
  \draw[rueckgabe] ([xshift=-0.09cm]7.2,-3.0) -- node[above,font=\scriptsize,inner sep=1pt] {\lsgtext{\code{3}}} ([xshift=0.09cm]3.8,-3.0);
  \draw[rueckgabe] ([xshift=-0.09cm]3.8,-3.35) -- node[above,font=\scriptsize,inner sep=1pt] {\lsgtext{\code{3}}} (0.6,-3.35);
\end{scope}
\end{tikzpicture}\\
{\raggedright\textbf{c)} Die Aufrufe sind \lsg[5.25cm]{synchron}, weil \lsg[7.2cm]{jeder Aufrufer auf die Antwort wartet}.\par}
\end{aufgabe}

\begin{merke}[Rekursion]
Eine Methode heißt \lsg[2.1cm]{rekursiv}, wenn sie sich selbst aufruft. Ohne \lsg[3.6cm]{Abbruchbedingung} ruft sie sich endlos auf. Die Liste ist eine \textbf{rekursive Datenstruktur}: Ein Knoten mit seinem Rest ist wieder eine Liste. Darum erledigt jeder Knoten seinen Teil und gibt den \lsg[2.2cm]{Rest} an seinen \lsg[2.4cm]{Nachfolger} weiter.
\end{merke}

\newpage
\begin{aufgabe}[Projekt Taxizentrale, TODO 3]{3}
\begin{minipage}[t]{0.52\linewidth}
\vspace{0pt}
\begin{java}[basicstyle=\ttfamily\scriptsize]
// Knoten: hängt d ans Ende an
public boolean add(DatenElement d) {
   if (§\lsg[2.5cm]{nachfolger == null}§) {
      nachfolger = §\lsg[2.6cm]{new Knoten(d)}§;
      return true;
   } else {
      return §\lsg[3.2cm]{nachfolger.add(d)}§;
   }
}
// Liste
public boolean add(DatenElement d) {
   if (anfang == null) {     // Sonderfall
      anfang = §\lsg[2.6cm]{new Knoten(d)}§;
      return true;
   } else {
      return §\lsg[2.7cm]{anfang.add(d)}§;
   }
}
\end{java}
\end{minipage}\hfill
\begin{minipage}[t]{0.45\linewidth}
\vspace{0pt}\raggedright
\textbf{a)} Ergänze \code{add} und übertrage es ins Projekt.\\[1.5mm]
\textbf{b)} Der letzte Knoten liefert \code{true}. Was geschieht mit diesem Wert auf dem Rückweg?\\
\lsg[0.98\linewidth]{Jeder Knoten gibt ihn mit return}\\[1mm]
\lsg[0.98\linewidth]{unverändert weiter bis zur Liste.}\\[1.5mm]
\textbf{c)} Warum braucht die Liste einen Sonderfall, der Knoten aber nicht?\\
\lsg[0.98\linewidth]{In der leeren Liste ist anfang null:}\\[1mm]
\lsg[0.98\linewidth]{Es gibt keinen Knoten, der fragen könnte.}
\end{minipage}
\end{aufgabe}

\begin{aufgabe}[Projekt Taxizentrale, TODO 4a bis 4c]{4}
\begin{minipage}[t]{0.47\linewidth}
\vspace{0pt}\raggedright
\textbf{a)} Gleich oder dasselbe? Sage die Ausgabe vorher:\\[-0.5mm]
\begin{java}[basicstyle=\ttfamily\scriptsize]
Taxi a = new Taxi(7, "Aylin");
Taxi b = new Taxi(7, "Aylin");
println(a == b);
\end{java}
\lsg[0.98\linewidth]{false: zwei verschiedene Objekte}\\[1mm]
\code{==} prüft, ob zwei Verweise auf \lsg[3.2cm]{dasselbe Objekt} zeigen. Für \textbf{inhaltliche} Gleichheit erweitern wir den Vertrag \code{DatenElement} (4a):\\[1mm]
\code{boolean} \lsg[3.4cm]{istGleich(DatenElement d)}\code{;}\\[1mm]
\textbf{b)} Danach meldet der Compiler einen Fehler in \code{Taxi}. Warum? Ergänze dort (4b):
\end{minipage}\hfill
\begin{minipage}[t]{0.5\linewidth}
\vspace{0pt}
\begin{java}[basicstyle=\ttfamily\scriptsize]
// Taxi (4b)
public boolean istGleich(DatenElement d) {
   return §\lsg[4.7cm]{gibSchluesselwert().equals(}§
          §\lsg[3.6cm]{d.gibSchluesselwert())}§;
}
// Knoten (4c): Ist d ab hier enthalten?
public boolean enthaelt(DatenElement d) {
   if (§\lsg[2.6cm]{daten.istGleich(d)}§) {
      return true;
   } else if (§\lsg[2.4cm]{nachfolger == null}§) {
      return §\lsg[0.9cm]{false}§;
   } else {
      return §\lsg[3.3cm]{nachfolger.enthaelt(d)}§;
   }
}
\end{java}
\lsg[0.98\linewidth]{b) Taxi erfüllt den erweiterten Vertrag noch nicht.}
\end{minipage}
\end{aufgabe}

\begin{aufgabe}[Projekt Taxizentrale, TODO 5]{5}
\begin{minipage}[t]{0.42\linewidth}
\vspace{0pt}\raggedright
Taxi 7 meldet sich ab. Ein Knoten kann sich nicht selbst aushängen: Den Verweis auf ihn hat sein \lsg[2.3cm]{Vorgänger}. Deshalb schaut jeder Knoten auf seinen \textbf{Nachfolger} voraus.\\[1mm]
\begin{tikzpicture}[node distance=5mm]
\node[obj] (a) {\textbf{k1}\nodepart{second}Taxi 12};
\node[obj,right=of a] (b) {\textbf{k2}\nodepart{second}Taxi 7};
\node[obj,right=of b] (c) {\textbf{k3}\nodepart{second}Taxi 3};
\draw[zeiger,gray,densely dashed] (a) -- (b);
\draw[gray] ($(a.east)!0.5!(b.west)+(-1mm,-1mm)$) -- ++(2mm,2mm);
\draw[zeiger] (b) -- (c);
\draw[zeiger,teal!60!black] (a.south) to[bend right=35] node[schritt,below=0.5mm] {2} (c.south);
\node[schritt,above=1mm of b] {1};
\end{tikzpicture}\\[1mm]
\scriptsize \nr{1} k1 prüft die Daten seines Nachfolgers.\\
\nr{2} Passen sie, überspringt k1 ihn.
\end{minipage}\hfill
\begin{minipage}[t]{0.55\linewidth}
\vspace{0pt}
\begin{java}[basicstyle=\ttfamily\scriptsize]
// Knoten
public boolean entfernen(DatenElement d) {
   if (nachfolger == null) {
      return false;
   } else if (§\lsg[4.6cm]{nachfolger.getDaten().istGleich(d)}§) { // §\nr{1}§
      nachfolger = §\lsg[3.8cm]{nachfolger.getNachfolger()}§; // §\nr{2}§
      return true;
   } else {
      return §\lsg[3.4cm]{nachfolger.entfernen(d)}§;
   }
}
\end{java}
\small Den ersten Knoten kann niemand überspringen. Sonderfall in \code{Liste.entfernen}:\\[0.5mm]
\code{if (anfang.getDaten().istGleich(d))}\\ \code{\ \ \ anfang = }\lsg[3.8cm]{anfang.getNachfolger()}\code{;}
\end{minipage}
\end{aufgabe}

\begin{hinweis}\small
\textbf{Es geht auch anders -- besprecht mündlich, warum man es so machen kann:}
\begin{itemize}
\item \code{laenge()} ginge auch mit einer Wiederholung wie \code{anzahl()}. Was ist leichter zu lesen, was braucht weniger Speicher?
\item Die Liste könnte wie die Warteschlange einen Verweis auf das \code{ende} haben. Was würde \code{add} dann sparen?
\item \code{entfernen} im Knoten könnte statt \code{boolean} den Knoten zurückgeben, der danach an seiner Stelle steht (sich selbst oder seinen Nachfolger). Was müsste der Aufrufer mit der Rückgabe tun?
\end{itemize}
\end{hinweis}

\end{document}
